% APA author-date reference list, rendered manually rather than through a .bst
% file. biber is unavailable in the build environment, apalike is not APA 7, and
% a hand-rendered thebibliography compiles wherever Overleaf puts it. natbib's
% authoryear mode takes the in-text form from the bracketed label.
%
% One deliberate departure from APA 7: the bracketed label supplies both the
% narrative and the parenthetical citation, so works with two authors read
% "Ledoit and Wolf (2003)" narratively and "(Ledoit and Wolf, 2003)" in
% parentheses, where APA would use an ampersand in the second. Labels are set
% for the narrative form because the text uses it predominantly. Reference-list
% entries themselves follow APA 7 exactly, ampersand included.
%
% Metadata mirrors references/references.bib, where every DOI was checked
% against Crossref. Do not edit an entry here without editing that file.
\begin{thebibliography}{99}

\bibitem[Alonso and Nusinzon(2020)]{alonso2020devil}
Alonso, N., \& Nusinzon, O. (2020). The devil is in the details: The risks
often ignored in low-volatility investing. \textit{The Journal of Portfolio
Management}, \textit{46}(7), 58--70.
\url{https://doi.org/10.3905/jpm.2020.1.163}

\bibitem[Ang et~al.(2006)]{ang2006crosssection}
Ang, A., Hodrick, R. J., Xing, Y., \& Zhang, X. (2006). The cross-section of
volatility and expected returns. \textit{The Journal of Finance},
\textit{61}(1), 259--299.
\url{https://doi.org/10.1111/j.1540-6261.2006.00836.x}

\bibitem[AQR Capital Management(n.d.)]{aqrbab}
AQR Capital Management. (n.d.). \textit{Betting against beta: Equity factors,
daily} [Data set]. Retrieved October 10, 2026, from
\url{https://www.aqr.com/Insights/Datasets/Betting-Against-Beta-Equity-Factors-Daily}

\bibitem[Bai and Ng(2002)]{bai2002determining}
Bai, J., \& Ng, S. (2002). Determining the number of factors in approximate
factor models. \textit{Econometrica}, \textit{70}(1), 191--221.
\url{https://doi.org/10.1111/1468-0262.00273}

\bibitem[Baker et~al.(2011)]{baker2011benchmarks}
Baker, M., Bradley, B., \& Wurgler, J. (2011). Benchmarks as limits to
arbitrage: Understanding the low-volatility anomaly. \textit{Financial Analysts
Journal}, \textit{67}(1), 40--54. \url{https://doi.org/10.2469/faj.v67.n1.4}

\bibitem[Blitz and van Vliet(2007)]{blitz2007volatility}
Blitz, D. C., \& van Vliet, P. (2007). The volatility effect. \textit{The
Journal of Portfolio Management}, \textit{34}(1), 102--113.
\url{https://doi.org/10.3905/jpm.2007.698039}

\bibitem[Blume(1971)]{blume1971assessment}
Blume, M. E. (1971). On the assessment of risk. \textit{The Journal of
Finance}, \textit{26}(1), 1--10.
\url{https://doi.org/10.1111/j.1540-6261.1971.tb00584.x}

\bibitem[Briner and Connor(2008)]{briner2008howmuch}
Briner, B., \& Connor, G. (2008). How much structure is best? A comparison of
market model, factor model and unstructured equity covariance matrices.
\textit{The Journal of Risk}, \textit{10}(4), 3--30.
\url{https://doi.org/10.21314/JOR.2008.180}

\bibitem[Clarke et~al.(2006)]{clarke2006minimum}
Clarke, R. G., de Silva, H., \& Thorley, S. (2006). Minimum-variance portfolios
in the U.S. equity market. \textit{The Journal of Portfolio Management},
\textit{33}(1), 10--24. \url{https://doi.org/10.3905/jpm.2006.661366}

\bibitem[DeMiguel et~al.(2009)]{demiguel2009optimal}
DeMiguel, V., Garlappi, L., \& Uppal, R. (2009). Optimal versus naive
diversification: How inefficient is the 1/N portfolio strategy? \textit{The
Review of Financial Studies}, \textit{22}(5), 1915--1953.
\url{https://doi.org/10.1093/rfs/hhm075}

\bibitem[Dom et~al.(2025)]{dom2025beyond}
Dom, M. S., Howard, C., Jansen, M., \& Lohre, H. (2025). Beyond GMV: The
relevance of covariance matrix estimation for risk-based portfolio
construction. \textit{Quantitative Finance}, \textit{25}(3), 403--419.
\url{https://doi.org/10.1080/14697688.2025.2468268}

\bibitem[Fama and French(2015)]{fama2015five}
Fama, E. F., \& French, K. R. (2015). A five-factor asset pricing model.
\textit{Journal of Financial Economics}, \textit{116}(1), 1--22.
\url{https://doi.org/10.1016/j.jfineco.2014.10.010}

\bibitem[Frazzini and Pedersen(2014)]{frazzini2014betting}
Frazzini, A., \& Pedersen, L. H. (2014). Betting against beta. \textit{Journal
of Financial Economics}, \textit{111}(1), 1--25.
\url{https://doi.org/10.1016/j.jfineco.2013.10.005}

\bibitem[French(n.d.)]{frenchdatalibrary}
French, K. R. (n.d.). \textit{Fama/French data library} [Data set]. Tuck School
of Business, Dartmouth College. Retrieved October 10, 2026, from
\url{https://mba.tuck.dartmouth.edu/pages/faculty/ken.french/data_library.html}

\bibitem[Goldberg et~al.(2022)]{goldberg2022dispersion}
Goldberg, L. R., Papanicolaou, A., \& Shkolnik, A. (2022). The dispersion bias.
\textit{SIAM Journal on Financial Mathematics}, \textit{13}(2), 521--550.
\url{https://doi.org/10.1137/21M144058X}

\bibitem[Haugen and Baker(1991)]{haugen1991efficient}
Haugen, R. A., \& Baker, N. L. (1991). The efficient market inefficiency of
capitalization-weighted stock portfolios. \textit{The Journal of Portfolio
Management}, \textit{17}(3), 35--40.
\url{https://doi.org/10.3905/jpm.1991.409335}

\bibitem[Invesco Exchange-Traded Fund Trust II(2024)]{invesco2024splv497k}
Invesco Exchange-Traded Fund Trust II. (2024). \textit{Invesco S\&P 500 Low
Volatility ETF summary prospectus} (Form 497K). U.S. Securities and Exchange
Commission.
\url{https://www.sec.gov/Archives/edgar/data/1378872/000119312524282757/d861761d497k.htm}

\bibitem[iShares Trust(2011)]{ishares2011usmv497k}
iShares Trust. (2011). \textit{iShares MSCI USA Minimum Volatility Index Fund
summary prospectus} (Form 497K). U.S. Securities and Exchange Commission.
\url{https://www.sec.gov/Archives/edgar/data/1100663/000119312511266729/d240179d497k.htm}

\bibitem[Jagannathan and Ma(2003)]{jagannathan2003risk}
Jagannathan, R., \& Ma, T. (2003). Risk reduction in large portfolios: Why
imposing the wrong constraints helps. \textit{The Journal of Finance},
\textit{58}(4), 1651--1683. \url{https://doi.org/10.1111/1540-6261.00580}

\bibitem[Ledoit and Wolf(2003)]{ledoit2003improved}
Ledoit, O., \& Wolf, M. (2003). Improved estimation of the covariance matrix of
stock returns with an application to portfolio selection. \textit{Journal of
Empirical Finance}, \textit{10}(5), 603--621.
\url{https://doi.org/10.1016/S0927-5398(03)00007-0}

\bibitem[Ledoit and Wolf(2004)]{ledoit2004honey}
Ledoit, O., \& Wolf, M. (2004). Honey, I shrunk the sample covariance matrix.
\textit{The Journal of Portfolio Management}, \textit{30}(4), 110--119.
\url{https://doi.org/10.3905/jpm.2004.110}

\bibitem[Markowitz(1952)]{markowitz1952portfolio}
Markowitz, H. (1952). Portfolio selection. \textit{The Journal of Finance},
\textit{7}(1), 77--91.
\url{https://doi.org/10.1111/j.1540-6261.1952.tb01525.x}

\bibitem[MSCI Inc.(2012)]{msci2012minvol}
MSCI Inc. (2012). \textit{MSCI minimum volatility indices methodology}.
Retrieved October 10, 2026, from
\url{https://www.msci.com/eqb/methodology/meth_docs/MSCI_Minimum_Volatility_Methodology_Jan12.pdf}

\bibitem[MSCI Inc.(2025)]{msci2025minvoladvanced}
MSCI Inc. (2025). \textit{MSCI minimum volatility advanced indexes
methodology}. Retrieved October 10, 2026, from
\url{https://www.msci.com/indexes/documents/methodology/3_MSCI_Minimum_Volatility_Advanced_Indexes_Methodology_20250203.pdf}

\bibitem[Newey and West(1987)]{newey1987simple}
Newey, W. K., \& West, K. D. (1987). A simple, positive semi-definite,
heteroskedasticity and autocorrelation consistent covariance matrix.
\textit{Econometrica}, \textit{55}(3), 703--708.
\url{https://doi.org/10.2307/1913610}

\bibitem[S\&P Dow Jones Indices(n.d.)]{spdji_lowvol_methodology}
S\&P Dow Jones Indices. (n.d.). \textit{S\&P 500 Low Volatility Index
methodology}.
\url{https://www.spglobal.com/spdji/en/documents/methodologies/methodology-sp-low-volatility-indices.pdf}

\bibitem[U.S. Securities and Exchange Commission(n.d.)]{secedgarapi}
U.S. Securities and Exchange Commission. (n.d.). \textit{EDGAR application
programming interfaces}. Retrieved October 10, 2026, from
\url{https://www.sec.gov/search-filings/edgar-application-programming-interfaces}

\end{thebibliography}
